%% Abstract + keywords (SPIE style).
%% Rev. 2 --- adapted from the pre-study draft to what the results actually show:
%% the RL policy does NOT dominate; it ties two very different architectures in the
%% realistic deployment condition. Every number is traceable to
%% papers/paper1-benchmark/README.md and the reports/ notebook.

\begin{abstract}
Micro-unmanned aerial vehicles operating outside the laboratory face sustained and
unsteady wind, payload changes that are a large fraction of their own weight, and
onboard state estimation far poorer than motion capture --- all on a platform with
little control authority to spare. Controllers for this regime are published on
different vehicles, references, metrics and sensing stacks, so their relative merits,
and especially their failure modes, cannot be compared across papers. We benchmark six
trajectory-tracking controllers, one per control paradigm, on a single simulated
brushless Crazyflie~2.1 ($43.4$~g, thrust-to-weight $1.88$), a single reference family
(figure-eight Lissajous at three cycle times) and a single command interface: a PID with
acceleration feedforward, active disturbance rejection control with an extended state
observer, sampling-based MPPI with $\mathcal{L}_1$ adaptation, an offset-free nonlinear
MPC, a hybrid stack pairing the observer with a pretrained learned adaptive rate
controller, and a reinforcement-learning policy trained with a privileged critic and
sensor-noise domain randomization. All six are evaluated under constant wind, gusts,
in-ground effect, a $23\,\%$ payload step, and a Lighthouse measurement model grounded
in published characterization data, using ten evaluation seeds and three training seeds
so that every claim carries an error bar. No controller dominates: the per-condition
champions differ, and each wins through a different disturbance-estimation mechanism.
Under the realistic deployment condition --- degraded onboard sensing with wind --- four
of them become statistically indistinguishable: the offset-free MPC at
$0.057\pm0.014$~m, the learned policy at $0.059\pm0.002$~m, the hybrid adaptive stack at
$0.060\pm0.009$~m and the extended state observer alone at $0.063\pm0.006$~m. Four
unrelated mechanisms saturating at the same value argues that the regime is
estimator-limited rather than architecture-limited. We
distil the study into seven transferable lessons, two of which overturn conclusions we
had drawn ourselves from single-seed, full-window data: mean-only tables hide leaders
that fail on unlucky sensor-bias draws, and an apparent noise fragility proved to be a
launch transient generic to the predictive family. Finally, we describe the hardware
validation protocol for the recommended controller --- a programmable fan-array wind
tunnel with motion-capture ground truth --- designed so that the physical disturbance
conditions correspond one-to-one with the simulated ones.
% TODO (after the wind-tunnel campaign): replace the last sentence with the measured
% sim-to-real outcome, e.g. "...; in flight, the recommended stack tracked to XX m
% under a YY m/s programmed gust field, an increase of ZZx over its simulated value."
\end{abstract}

\keywords{quadrotor control, trajectory tracking, benchmarking, disturbance rejection,
active disturbance rejection control, model predictive control, reinforcement learning,
domain randomization, wind tunnel, motion capture, nano-aerial vehicles}
